\subsection{Four matching witnesses over a quadratic circuit with excess one}
Full-graph safety provides the missing restriction on insertion contexts when the quadratic circuit is larger than its minimum. We first isolate the functional equality case in the four-product argument.

\begin{lemma}[The four-product matching alternative]\label{lem:fourproductalternative}
Suppose a positive circuit has nonempty squarefree quadratic support, six deleted gates consisting of four products and two additions, and eight vertices on which its Boolean output is a four-edge matching while its quadratic part vanishes. After full-carrier erasure, the only possible restricted computation has two visible cuts, zero independent-symbol pair coefficient, and the form
\[
 gL\vee fM,
\]
where $g$ and $f$ are independent products of three old linear clauses. Each safe quartic witnesses exactly two matching edges, and these two pairs are disjoint. Its four restricted clauses are exactly the four two-element transversals of its witnessed pair.
\end{lemma}
\begin{proof}
The following is the functional portion of \citet[proof of Lemma~5.35]{Companion}, before its use of quadratic optimality. We give the argument to specify the equality case needed here. Full-carrier erasure cannot remove a private gate. Otherwise the exact quadratic extraction would be preserved with at most five deletions. The four-product obstruction still forces at least four private products, and a missing witness forces at least one visible cut. The only possible profile would then have four products and one addition. Its unique high input is a multiplication DAG with at most five old linear boundary clauses, and the output context adds at most one clause. The resulting matching would have a CNF with at most six clauses, contrary to the required sixteen. This essentiality concerns the original circuit, before restriction.

If there is just one visible cut, all private gates feed its high operand. It has four products and at most one private addition over old linear clauses. Lemma~\ref{lem:fourandtwoor}, followed by the affine direct context, gives at most thirteen CNF clauses. The matching requires sixteen, so both additions are visible cuts. Their high inputs are private products $g,f$, and the exact restricted frame is
\[
 F=gL\vee fM\vee\delta gf,\qquad\delta\in\{0,1\}.
\]
The contexts are arbitrary affine positive clauses at this stage. Let $k$ count the private products ancestral to the first high $g$ in the original erased circuit. Every other private product feeds $f$. Keep the four nodes even if some restricted values vanish or coincide.

If $k=4$, the second high is an ancestor of $g$ or equals it. Write $g=fA$ Booleanly. Their total number of distinct old factors is at most five, and
\[
 F=f(AL\vee M\vee\delta A)
\]
has at most six clauses after distribution and absorption.

Suppose $k=3$. Then $g$ has at most $b=4$ old factors and only one product remains for $f$. If both inputs of $f$ are private ancestors of $g$, its factors are among those of $g$, so the preceding factorization applies. If exactly one input is a proper private ancestor $u$, write $g=uW$. The other input gives $f=uA$ or $uA\vee g$, including zero old parts. If $u$ has $r\ge1$ factors and $W$ has $t$ remaining factors, then $r+t\le4$. Factoring $u$ leaves, respectively,
\[
 WL\vee AM\vee\delta WA,
 \qquad W(L\vee M\vee\delta)\vee AM.
\]
The bracket has at most $2(t+1)$ clauses; with $u$ this gives at most $r+2(t+1)\le9$. For the first bracket with $\delta=1$, the identity
\[
 WL\vee AM\vee WA=(W\vee A)(W\vee M)(A\vee L)
\]
gives the slightly smaller count $2t+1$.

Without a proper private-ancestor operand, $f$ is $AB$, $AB\vee gB$, or $AB\vee g$. Direct distribution gives the following bounds:
\[
\begin{array}{c|cc}
f&\delta=0&\delta=1\\\hline
AB&3(b+1)\le15&3b+2\le14\\
AB\vee gB&3b+2\le14&3b+1\le13\\
AB\vee g&3b+1\le13&3b\le12.
\end{array}
\]
For example, in the middle row with $\delta=0$ the identity is
$F=(g\vee ABM)(L\vee BM)$; with $\delta=1$ it is
$F=(g\vee ABM)(L\vee B)$. In the last row absorption gives
$F=(g\vee ABM)(L\vee M)$ or $F=g\vee ABM$.
All these counts are below sixteen.

Suppose $k=2$. Write $u=UV$ and $g=uW$, allowing $W=1$ when the second product repeats its input. If the two-product network for $f$ reads $u$ directly, that factor persists to $f$. Under it a $g$-affected boundary form $A\vee g$ becomes $A\vee W$. Thus $f=uK$, with $K$ a product of at most two linear clauses, and
\[
 F=u(WL\vee KM\vee\delta WK).
\]
The bracket has at most six clauses for $\delta=0$, and at most five for $\delta=1$ by
$(W\vee K)(W\vee M)(K\vee L)$. Including $U,V$ gives at most eight.

Otherwise $g$ is a product of at most $b=3$ old clauses, and $f$ is a product of at most $t=3$ old or $g$-affected clauses. This leaf count includes sharing: a connected two-product multiplication graph has at most three distinct boundary inputs. If none is affected, write $f=K$. For $\delta=0$ the clause count is at most $(b+1)(t+1)\le16$; for $\delta=1$ the identity
\[
 gL\vee KM\vee gK=(g\vee K)(g\vee M)(K\vee L)
\]
gives $bt+b+t\le15$ clauses.

If at least one factor is affected, distributivity and absorption write
$f=H_0K_0\vee gH_0$, where $H_0$ has at most $h\le2$ old factors and $H_0K_0$ has at most $t\le3$. Then
\[
 F=(g\vee H_0K_0M)(L\vee H_0M)\quad(\delta=0),
\]
and the second factor becomes $L\vee H_0$ for $\delta=1$. The counts are at most $b(t+1)+h+1\le15$ and $b(t+1)+h\le14$.

Finally, if $k=1$, $g$ has at most $b=2$ factors and the remaining three products give at most $t=4$ boundary factors for $f$. If none is affected, the preceding counts are at most fifteen and fourteen. Otherwise $h\le3$ and the counts are at most $b(t+1)+h+1\le14$ and $b(t+1)+h\le13$. A directly used $g$ is included by setting the old part of its affected clause to zero. Repeated factors and restricted zero values only decrease these bounds.

The only count that can reach sixteen is therefore $k=2$, $\delta=0$, with $g,f$ independent products of exactly three old clauses. Each summand $gL$ or $fM$ is safe and has at most four clauses. A safe CNF true on $s$ matching edges needs at least $2^s$ clauses, by the independent-transversal argument; hence each summand witnesses at most two. They must each witness two, with disjoint pairs, to cover all four edges.

For a safe quartic witnessing two edges, each clause meets both edge assignments. The four independent transversals must all be rejected, and one clause can reject at most one. Equality therefore forces the four distinct complementary two-element clauses. No clause can contain an outside variable: adjoining that variable to its uniquely rejected transversal would make all four clauses true on an independent assignment. A clause cannot be constant either. These are exactly the four asserted transversals.
\end{proof}

\begin{proposition}[Full insertion contexts force a small cover]\label{prop:sixmatchingcover}
Let $G$ have no four-cycle, and let a positive variable-input binary circuit compute $P\in\Sup(G)$ with nonempty quadratic support $H$. Suppose exactly four edges lack quadratic witnesses and their eight endpoints induce precisely those four edges in $G$. If exactly six gates are deleted, four of them products, and $K$ products remain in the quadratic circuit, then
\[
 \tau(G)\le K+2.
\]
Consequently this circuit has at least $|E(G)|+\tau(G)-1$ gates.
\end{proposition}
\begin{proof}
Apply full-carrier erasure and restrict to the eight endpoints. By Lemma~\ref{lem:fourproductalternative}, the sole possible computation has two visible cuts and the independent-symbol form $P_2+gD_g+fD_f$, with zero pair coefficient. Here $D_g,D_f$ are the \emph{full} old direct contexts, before restriction. Relabel the missing edges
\[
 aa',\quad bb',\quad cc',\quad dd'
\]
so that the restricted contexts are $a\vee b$ and $c\vee d$. Quartic equality gives
\[
 g(\{a',b'\})=1,\qquad f(\{c',d'\})=1.
\]
All these statements concern Boolean supports; positive coefficients do not change them.

Let $A_g$ be the nonzero old linear low part of the $g$ cut. The product $A_gD_g$ occurs in the squarefree quadratic frame. Since $D_g$ contains two distinct variables $a,b$, four-cycle-freeness forces $A_g$ to have singleton support, say $\{h\}$: two distinct variables in each disjoint factor would give a four-cycle. Squarefreeness makes the factor supports disjoint. The edges $ha,hb$ belong to $H$, so $h$ lies outside the eight endpoints, which contain no $H$-edge.

Suppose $D_g$ contains an additional vertex $z$. Its restriction has only $a,b$, so $z$ also lies outside the eight endpoints. The edge $hz$ lies in $H$, and $z\ne h$ by squarefreeness. On the assignment $\{a',b',z\}$ the positive contribution $gD_g$ is true. Since $a'b'$ is not an edge, correctness forces $za'$ or $zb'$ to be an edge of $G$. In the first case $h,a,a',z$ form a four-cycle, and in the second $h,b,b',z$ do. All four vertices are distinct. This contradiction proves
\[
 \operatorname{supp}(D_g)=\{a,b\}.
\]
The identical argument gives $\operatorname{supp}(D_f)=\{c,d\}$. The full contexts, not merely their restrictions, are disjoint two-element sets.

Call a quadratic product affected by a cut if its independent degree-one symbol reaches one of the product's two linear factors. Each cut affects at least one product, because its direct context is nonzero and linear. It cannot affect both factors of one product: substituting its nonzero old low part would produce a square. If a product is affected by the $g$ cut, its old opposite factor contributes positively to $D_g$, so its support is contained in $\{a,b\}$. Thus these two vertices cover every edge of that product star, including edges from other old summands in the perturbed factor.

No product is affected by both cuts. If both symbols enter the same factor, the nonempty opposite factor contributes to both disjoint contexts. If they enter opposite factors, they produce a positive independent-symbol pair coefficient, which has been excluded. Both conclusions use positivity, so there is no cancellation. They allow shared inputs, shared intermediate values and nested cut locations.

The two disjoint affected sets therefore contain at least two quadratic products. Their edges are covered by $S=\{a,b,c,d\}$, which also covers all four missing edges. For each unaffected product choose one star center. There are at most $K-2$ such products, and their centers together with $S$ cover $G$. Hence $\tau(G)\le K-2+4=K+2$.

We justify the intermediate count $|C_2|\ge |E(H)|+K-1$ without an optimality assumption. Every quadratic product is a star: both factors are linear, their supports are disjoint by squarefreeness, and two variables in each factor would create a four-cycle. Choose its singleton factor as its center, choosing either endpoint for a one-edge product. Suppress linear additions with singleton support; this preserves all $K$ products and the positive edge support, and only decreases size.

A remaining linear wire has at least two variables in its support. It cannot feed products with distinct centers, since those two centers and any two of its variables would form a four-cycle. Thus these linear gates partition by their unique product center. For a center $v$, let $k_v$ be the number of its products, $q_v$ the number of distinct variables in the union of their leaf factors, and $a_v$ the number of linear additions assigned to it. The binary input graph for these factors has at most $k_v$ connected components, each containing one of the factor outputs. Counting its vertices and edges gives $a_v\ge q_v-k_v$. This also includes single-variable factor outputs and repeated use of one factor wire. Every edge of $H$ occurs in at least one chosen star, so $\sum_vq_v\ge|E(H)|$ and $\sum_vk_v=K$. The linear part therefore has at least $|E(H)|-K$ additions. Above the $K$ distinct product nodes the output network consists only of additions and has at least $K-1$ gates by the same connected-graph count. Adding the $K$ products proves the claimed inequality.

Since $|E(G)|=|E(H)|+4$ and six gates were deleted,
\[
 |C|\ge |E(H)|+K+5
       =|E(G)|+K+1
       \ge |E(G)|+\tau(G)-1.
\]
\end{proof}

\begin{theorem}[Maximal cover gain with four missing witnesses]\label{thm:fourmaxgain}
Let $G$ be a finite simple graph without a four-cycle. If a positive variable-input binary circuit computes a polynomial in $\Sup(G)$ with exactly four missing quadratic edge witnesses and nonempty quadratic support $H$, and if $\tau(G)-\tau(H)=4$, then its size is at least $|E(G)|+\tau(G)-1$.
Thus the remaining possible pairs for a four-witness counterexample are
\[
\begin{array}{c|l}
d&(D,e)\\\hline
2&(5,0)\\
3&(5,0),(5,1),(6,0).
\end{array}
\]
\end{theorem}
\begin{proof}
The previous reductions leave only $(D,e)=(6,1)$ when $d=4$. Maximal cover gain forces the missing edges to be an induced four-edge matching with no old edge on their endpoints. Exactly four products and two additions are deleted. Proposition~\ref{prop:sixmatchingcover} excludes a saving in this final case. The displayed smaller-gain alternatives are those of the preceding reduction.
\end{proof}
